\section{Линейно-разделимый и линейно-неразделимый случаи}

Ранее мы уже успели познакомиться с этими понятиями, и самое время рассказать зачем.

Как мы уже поняли, задача SVM заключается в том, чтобы попытаться максимизировать расстояние от разделяющей гиперплоскости до точек каждого из классов. Но тут возникает закономерный вопрос "--- а что делать при линейно неразделимом случае?

И тут мы плавно переходим к факту, что на самом деле, стратегии для обоих этих случаев разная, и далее мы будем выводить их для линейно разделимого и линейно неразделимого случаев.

\subsection{Линейно-разделимый случай}

Пусть у нас есть обучающая выборка
\[
\mathcal{D} = \{(x_i, y_i)\}_{i=1}^n, \quad y_i \in \{-1, +1\}.
\]
Мы говорим, что эта выборка \textit{линейно разделима}, если существует такая гиперплоскость
\[
w^\top x + b = 0,
\]
которая полностью отделяет точки одного класса от точек другого, то есть выполняются неравенства
\[
y_i (w^\top x_i + b) > 0 \quad \forall i = 1, \dots, n.
\]

Среди всех гиперплоскостей, удовлетворяющих этому условию, метод опорных векторов выбирает ту, для которой норма вектора весов минимальна:
\[
\max_{w} \frac{2}{||w||} 
\quad \text{при условии} \quad 
y_i (w^\top x_i + b) \ge 1, \; \forall i.
\]

Говоря проще "--- мы хотим максимизировать воображаемую ширину разделяющей гиперплоскости при условии, что она по прежнему корректно классифицирует все объекты.

Давайте в одну систему соберём все нужные условия:
\[
\begin{cases}
    $\|w\| \rightarrow \min$ \\
    $y_i (\langle w,x\rangle + b) \ge 0$ \text{Требование корректной классификации} \\
    $\vert \langle w, x\rangle + b\vert \ge 1$ \text{Ранее поставленное требование}
\end{cases}
\]

Давайте попробуем упростить эту систему. Заметим, что второе и третье условия уж слишком похожи. Мы можем подставить третье во второе и сказать, что всё это не меньше единицы, т.к. $y_i \in \{-1, 1\}$, и в таком случае наше требование корректной классификации не нарушится. 
\[
y_i(\langle w, x\rangle + b) \ge 1
\]

Вставим это в систему:
\[
\begin{cases}
    $\|w\| \rightarrow \min$ \\
    $y_i (\langle w,x\rangle + b) \ge 1$ \text{Требование корректной классификации} 
\end{cases}
\]

Так формулируется решение задачи с использованием \textbf{жёсткого зазора (hard margin SVM)} — базовый вариант метода, применимый к идеально разделимым данным.

\subsection{Линейно-неразделимый случай}

К сожалению, реальность редко бывает столь благосклонной.  
В большинстве реальных задач данные оказываются \textit{линейно неразделимыми}: классы частично пересекаются, в выборке есть шум, выбросы и т.д.

В этом случае не существует пары параметров $(w, b)$, для которых все точки удовлетворяли бы неравенству $y_i (w^\top x_i + b) > 0$.  
Сколько бы мы ни пытались повернуть или сдвинуть гиперплоскость, всегда найдутся точки, которые окажутся «по неправильную сторону».

Идея заключается в том, чтобы разрешить некоторым точкам нарушать ограничение $y_i (w^\top x_i + b) \ge 1$, но ввести за это «штраф».  
Для этого в модель добавляются специальные неотрицательные переменные $\xi_i$, называемые \textit{переменными послабления} (англ.~\textit{slack variables}).  
Каждая из них показывает, насколько сильно нарушено требование разделимости для соответствующей точки. Этими переменными мы намерены ослаблять наше <<модифицированное>> условие:
\[
y_i (w^\top x_i + b) \ge 1 - \xi_i, \quad \xi_i \ge 0.
\]

Если $\xi_i = 0$, точка классифицирована корректно и находится за пределами зазора.  
Если $0 < \xi_i < 1$, она попадает внутрь зазора, но остаётся на правильной стороне гиперплоскости.  
Если же $\xi_i > 1$, классификатор ошибся: точка оказалась не только в зазоре, но и с другой стороны разделяющей границы.

Однако у этого проблема есть одна, довольно неочевидная проблема: при $\xi \rightarrow \infty$ появляется решение $w = 0$. Разумеется, с такими весами модель ничему и не научится, поэтому нужно что"=то придумывать.

А придумаем мы вот что: в качестве условия добавим, что мы хотим минимизировать сумму $\xi$. В таком случае, в силу постановки задачи мы просто не сможем ставить сколь угодно большие ослабления условия.

Обычно сумма домножается на некоторый коэффициент $C > 0$, который задаёт \textit{цену ошибки}.  
\[
\max_{w, \xi} \frac{2}{\|w\|} + C \sum_{i=1}^n \xi_i
\quad \text{при условиях} \quad
y_i (w^\top x_i + b) \ge 1 - \xi_i, \quad \xi_i \ge 0.
\]

Теперь нам предстоит перейти к \textit{безусловному виду}, где эти условия будут подразумеваться самой формулой.

Для первого условия выразим $\xi$:
\[
\xi  \ge 1 - y_i(\langle w, x\rangle + b)
\]

Заметим, что в совокупности со вторым условием, $\xi$ должно быть не меньше нуля и числа $1-M$. Значит, оно равно
\[
\max (0, 1 -M)
\]

Теперь мы можем записать задачу в безусловном виде. Сразу запишем в привычном виде "--- функции потерь принято минимизировать, а я напомню, что мы выводим именно её:
\[
||w|| + C \frac{1}{\ell}\sum_{i=1}^\ell \max (0, 1-M) \rightarrow \min_w
\]

Таким образом выводится задача для \textbf{мягкого зазора (soft margin SVM)}.

